\textbf{APPENDICE}

A distanza di ventitré anni dalla prima pubblicazione in digitale (1992)
e di undici anni da quella cartacea (edizioni Phasar, Firenze, 2004),
penso sia opportuno aggiungere un paio di partite utili a sviluppare
alcuni concetti trattati in precedenza. Le due partite che presento,
giocate quest'anno nel Campionato Italiano a squadre, mi permettono,
nella prima partita, di spiegare meglio la forza e la debolezza del
pedone debole e di mostrare le possibilità offerte da un attacco
contemporaneo su entrambi i lati della scacchiera. La seconda partita
mostra come limitandosi a sviluppare in modo naturale i pezzi si possa
finire in vantaggio senza avere fatto praticamente niente.

\textbf{Gazzarri -- Leoncini,} Cecina 2015 (Difesa Siciliana)

\textbf{1.e4 c5 2.c3 d5 3.exd5 £xd5 4.d4 cxd4 5.cxd4} Sulla scacchiera
si è creato un pedone debole in d4. A differenze del vantaggio di
sviluppo che può essere riassorbito, la modifica strutturale degli
elementi presenti sulla scacchiera (struttura di pedoni e Alfieri)
possono rimanere per utta la partita e caratterizzarla dal momento in
cui compaiono fino alla fine. Talvolta la creazione di un pedone isolato
è frutto di manovre avversarie. In questo caso il Bianco ha deciso di
crearsene uno. Il pedone d4 costituisce per tutto il centro partita un
vantaggio perché permette il controllo delle case centrali e il rapido
spostamento dei pezzi. Se il Bianco non riesce a valorizzare questo
vantaggio prima del finale rischia di finale, una volta cambiati i pezzi
leggeri e con le sole torri in gioco, si troverà a difendere il pedone e
al massimo potrà aspirare a faticoso pareggio. Questo significa che un
pedone isolato indica quale sarà la strategia dei giocatori: il Bianco
dovrà rimanere nel centro partita più a lungo possibile mentre il Nero,
al contrario, potrà permettersi di cambiare i pezzi leggeri e avviarsi
al finale. \textbf{5\ldots ¤c6 6.¤f3 e5} Abbiamo già visto in precedenza
questa mossa (Corvi-Leoncini) che stravolge i binari su cui si stava
dirigendo la partita. Il Nero non accetta il gioco proposto dal bianco e
con un cambio elimina il pedone isolato bianco. \textbf{7.¤c3 ¥b4 8.¥d2
¥xc3 9.¥xc3 e4 10.¤e5 ¤xe5 11.dxe5 ¤e7 12.£e2} Le alternative sono
12.£a4+ e 12.£c2. Questa è una variante minore ma non meno velenosa.
\textbf{12\ldots{} 0--0 13.¦d1 £xa2 14.£xe4 ¥f5 15.£c4} Interessante, ma
non sembra dare più della parità, 15.£xb7 ¦ad8! 16.£xe7 (se 16.¥e2 ¤d5!)
16\ldots{} ¦xd1+ 17.¢xd1 £b1+ 18,¢d2 £c2+ 19.re3 £e4+ e patta per scacco
perpetuo. \textbf{15\ldots{} £xc4 16.¥xc4 ¦ac8 17.¥b3} Vale la pena di
fare il punto della situazione. Il vantaggio del Bianco risiede nella
coppia degli Alfieri, quello del Nero in due pedoni contro uno sul lato
di Donna, con i Re sul lato opposto. \textbf{17\ldots{} ¤c6 18.0--0 b5!}
Una mossa molto forte che specula sulla poca mobilità dell'¥c3 che non
ha buone case per ritirarsi. Il Nero gioca per vincere e per ottenere lo
scopo deve innanzitutto eliminare la coppia degli Alfieri al Bianco
cambiando un Cavallo con l'Alfiere bianco delle case nere (non con
l'altro altrimenti la partita, con gli Alfieri di colore contrario,
rischierebbe di finire patta). \textbf{19.¦fe1 b4 20.¥d4 ¤xd4} Scopo
raggiunto! Il Bianco non ha più i due Alfieri mentre al nero rimane il
due contro uno sul lato di Donna.\textbf{21.¦xd4 a5 22.g4?!} Gli Alfieri
rimasti sulla scacchiera sono entrambi sulle case bianche. Questo
significa che i due giocatori devono mettere i propri pedoni su quelle
nere. \textbf{22\ldots{} ¥e6 23.¥a4 g5!} Blocca il Pg4, possibile futuro
obiettivo dell'¥e5, su casa bianca. \textbf{24.¢g2 ¦c5} Interessante era
anche 24\ldots ¦fd8 25.¦xd8+ ¦xd8 26.¦d1!) \textbf{25.¢g3 ¦fc8 26.h4 h6
27.f4 gxf4+ 28.¢xf4 ¦c1 29.¦xc1 ¦xc1 30.g5 ¦f1+ 31.¢g3 hxg5 32.¦d8+ ¢g7
33.¦a8 gxh4+ 34.¢xh4}

La posizione merita un diagramma. Ora comincia un attacco portato
contemporaneamente sui due lati della scacchiera.

\includegraphics[width=2in,height=2in]{assets/appendix/media/image1.png}

\textbf{34\ldots{} ¦f2!} Il Nero attacca il Pb2, ma l'idea non è tanto,
o solo, quella di attaccare un pedone, altrimenti avrebbe raggiunto più
facilmente l'obiettivo con 34\ldots{} ¦b1 ma quella di creare una rete
di matto. \textbf{35.b3 ¦f3 Inseguendo in apparenza il pedone ma, in
caso di 36.¦xa5 replicare con 36\ldots.¢h6!! con matto imparabile alla
successiva con 37\ldots{} ¦h3\#.}

\textbf{Borselli -- Leoncini,} Firenze 2015 (Tartakower)

\textbf{1.d4 d5 2.c4 e6 3.¤c3 ¤f6 4.¥g5 ¥e7 5.e3 0--0 6.¤f3 h6 7.¥h4 b6
8.£b3} Una variante un pretenziosa. Il punto b6 è molto solitop e una
pressione sul pedone è destinata all'insuccesso. \textbf{8\ldots ¥b7} Il
nero continua, tranquillo, lo sviluppo. \textbf{9.c5} Coerente con la
mossa precedente ma sbagliata. \textbf{9\ldots ¤bd7 Il nero continua a
svilupparsi. 10.cxb6} Il mangiatore di tempi (c2-c4-c5-cxb6) finisce la
sua corsa dopo tre mosse rivelatesi completamente inutili.
\textbf{10\ldots{} axb6} e non colo pedone c perché la spinta tematica è
c7-c5. 11.¥e2 ¥a6 In modo da cambiare l'Alfiere cattivo. \textbf{12.¥xa6
¦xa6 13.0--0 c5 14.dxc5 bxc5 15.¤b5 ¤e4 16.¥xe7 £xe7 17.£c2 ¦b8} Tutte
le forze sono mobilitate. E il nero ha ottenuto questo senza il minimo
sforzo, limitandosi a sviluppare i pezzi. \textbf{18.a4 c4} Per dare
spazio a un Cavallo in c5 da dove attaccherà il Pa4 e controllerà le
case avanzate b3 e d3. \textbf{19.¤fd4 £b4!} Entra nel pollaio.
\textbf{20.¦fd1 ¤ec5 21.¦db1 ¦xa4!} Passa in vantaggio di materiale
senza avere fatto praticamente niente. La mossa più difficile (meno
naturale?) è stata 19\ldots £b4 ma le altre le poteva giocare chiunque
avesse deciso di basarsi su principi generali. \textbf{22.¤c6 £xb5
23.¤xb8 ¤xb8 24.¦xa4 £xa4 25.£xa4 ¤xa4} I due Cavalli valgono più della
Torre. Inoltre il Nero ha altri due vantaggi: un pedone in più che può
diventare facilmente libero e la possibilità di tenere chiuse le
colonne. \textbf{26.¦a1 ¤b6 27.b3 cxb3} Il tentativo bianco di aprire le
colonne è giusto ma i due Cavalli riesco a costruire un muro.
\textbf{28.¦b1 ¢f8 29.¦xb3 ¤8d7 30.¦b4 ¢e7 31.f3 e5 32.¢f2 ¢d6 33.h4 h5
34.g4 ¢c5 35.¦b1 ¤f6 36.g5 ¤g8!} Con l'idea di andare in f5 o g6, via
e7, per attaccare il pedone in h4. \textbf{37.¦a1 ¤e7 38.e4?} Affretta
la fine perché aiuta l'avversario a crearsi un pedone passato per giunta
sostenuto. \textbf{38\ldots{} d4 39.¦a5+ ¢d6 40.¦a6 ¢c6 41.¦a5 ¤g6
42.¢g3} Se 42.¦a7 ¤d7 \textbf{42\ldots{} ¤c4!} Inizio della manovra
vincente.

\includegraphics[width=2in,height=2in]{assets/appendix/media/image2.png}

\textbf{43.¦a8 d3 44.¢f2 ¤f4 45.¢e1 d2+ 46.¢d1 ¤d3 47.¦c8+} 47.¦d8 ¤d6
\textbf{47\ldots{} ¢b5 48.¦xc4 ¢xc4 0-1}
